\documentclass[12pt, a4paper]{article}

\usepackage[utf8]{inputenc}
\usepackage[T1]{fontenc}
\usepackage[brazil]{babel}
\usepackage{csquotes}
\usepackage[top=3cm, bottom=2cm, left=3cm, right=2cm]{geometry}
\usepackage{lmodern}
\usepackage{microtype}
\usepackage{amsmath, amssymb, mathtools}
\usepackage{booktabs}
\usepackage{multirow}
\usepackage{array}
\usepackage{tabularx}
\usepackage{graphicx}
\usepackage{float}
\usepackage{caption}
\usepackage{subcaption}

% ── Pacotes ────────────────────────────────────────────────────────────────────
\usepackage[utf8]{inputenc}
\usepackage[T1]{fontenc}
\usepackage[brazil]{babel}
\usepackage{amsmath, amssymb, amsthm}
\usepackage{mathtools}
\usepackage{geometry}
\usepackage{hyperref}
\usepackage{booktabs}
\usepackage{multirow}
\usepackage{array}
\usepackage{xcolor}
\usepackage{tcolorbox}
\usepackage{caption}
\usepackage{siunitx}   % alinhamento numérico nas tabelas
 
\geometry{margin=2.5cm}
 
% ── Cores e caixas ─────────────────────────────────────────────────────────────
\definecolor{azul}{RGB}{0,70,127}
\definecolor{fundo}{RGB}{240,247,255}
\tcbuselibrary{skins, breakable}
\newtcolorbox{modelbox}{
  colback=fundo, colframe=azul,
  fonttitle=\bfseries\color{white}, coltitle=white,
  title=Especificação do Modelo,
  breakable, enhanced
}
 
% ── Macros ─────────────────────────────────────────────────────────────────────
\newcommand{\T}{{\top}}
\newcommand{\mY}{\mathbf{Y}}
\newcommand{\mb}{\boldsymbol{\beta}}
\newcommand{\beps}{\boldsymbol{\varepsilon}}
\DeclareMathOperator{\GI}{GI}
\DeclareMathOperator{\Normal}{\mathcal{N}}
\newcommand{\sig}{\sigma^{2}_{\varepsilon}}


\usepackage{listings}
\lstset{
  backgroundcolor=\color{cinzaclaro},
  basicstyle=\ttfamily\small,
  breaklines=true,
  frame=single,
  language=R
}

\usepackage[colorlinks=true, linkcolor=azulescuro,
            citecolor=azulescuro, urlcolor=azulescuro]{hyperref}
\usepackage{fancyhdr}
\pagestyle{fancy}
\fancyhf{}
\fancyhead[L]{\small Regressão Linear com enfoque Bayesiano}
\fancyhead[R]{\small \thepage}
\renewcommand{\headrulewidth}{0.4pt}

\usepackage{setspace}
\onehalfspacing
\usepackage{enumitem}
\usepackage{parskip}

\usepackage{graphicx} % Required for inserting images

\begin{document}
% ------- CAPA -------
\begin{titlepage}
  \centering
  \vspace*{1cm}

  {\large Disciplina: Estatística Bayesiana}

  \vspace{2.5cm}

  \rule{\linewidth}{0.5pt}\\[0.4cm]
  {\LARGE \textbf{Regressão Linear com enfoque Bayesiano}\\[0.3cm]
  \large \textit{Combined Cycle Power Plant}}

  \rule{\linewidth}{0.5pt}

  \vspace{2cm}

  {\large \textbf{Autores:}}\\[0.3cm]
  {\large }\\[0.2cm]
  {\large }\\[0.2cm]
  {\large }\\[0.3cm]

  \vfill

  {\large SP}\\[0.2cm]
  {\large \today}
\end{titlepage}


% =================== SUMÁRIO ===================
\tableofcontents
\newpage

\section*{Introdução}
\addcontentsline{toc}{section}{Introdução}

Neste trabalho, será utilizada uma abordagem de regressão linear bayesiana para modelar a relação entre variáveis ambientais e a produção líquida horária de energia elétrica, medida em megawatts (MW). Para isso, emprega-se o conjunto de dados \textit{Combined Cycle Power Plant}, disponibilizado pela \textit{UCI Machine Learning Repository}. O banco contém (n=9568) observações coletadas ao longo de seis anos (2006--2011), durante períodos em que a usina operou sob carga total. As observações correspondem a médias horárias obtidas por sensores instalados no sistema de geração.

A variável resposta considerada neste estudo é a potência elétrica líquida produzida (\texttt{PE}), expressa em MW. Como variáveis explicativas são utilizadas quatro medidas ambientais contínuas: temperatura ambiente (\texttt{AT}), vácuo de exaustão (\texttt{V}), pressão atmosférica (\texttt{AP}) e umidade relativa do ar (\texttt{RH}). Essas variáveis representam condições externas que afetam diretamente o funcionamento das turbinas e, consequentemente, o desempenho energético do sistema.

Diferentemente da regressão linear sob a abordagem clássica, a inferência bayesiana trata os parâmetros do modelo como variáveis aleatórias e incorpora conhecimento prévio por meio de distribuições \textit{a priori}. Nesse contexto, o modelo é especificado pela combinação entre a função de verossimilhança e distribuições prévias para os parâmetros desconhecidos, resultando em distribuições posteriores obtidas pela aplicação do Teorema de Bayes.

Especificamente, será adotado o modelo de regressão linear gaussiano com prior conjugado \textit{Normal Gama Inversa} para o par $(\boldsymbol{\beta},\sigma^2)$. Essa escolha permite derivação analítica da distribuição posterior conjunta e fornece interpretação probabilística direta sobre os coeficientes de regressão e sobre a variabilidade residual do processo.

% Ao final, o modelo será utilizado para investigar como as condições ambientais estão associadas à produção de energia elétrica e para avaliar o desempenho preditivo da regressão linear sob uma perspectiva bayesiana.


\section*{Conceito teórico: Regressão Linear Bayesiana com Prior Normal–Gama Inversa}
\addcontentsline{toc}{section}{Conceito teórico: Regressão Linear Bayesiana com Prior Normal–Gama Inversa}

\subsection*{1. Modelo Generativo}
\addcontentsline{toc}{subsection}{1. Modelo Generativo}

Considere o modelo de regressão linear clássico:
\[ \mathbf{Y} = X\boldsymbol{\beta} + \boldsymbol{\varepsilon} \]

onde: $\mathbf{Y}\in\mathbb{R}^{n}$, $X\in\mathbb{R}^{n\times p}$, $\boldsymbol{\beta}\in\mathbb{R}^{p}$ e $\boldsymbol{\varepsilon}\sim\mathcal{N}_n(\mathbf{0},\sigma^2I_n)$.

Logo,
\[ \mathbf{Y}\mid\boldsymbol{\beta},\sigma^2 \sim \mathcal{N}_n(X\boldsymbol{\beta},\sigma^2I_n). \]

A densidade da verossimilhança é:
\[ p(\mathbf{Y}\mid\boldsymbol{\beta},\sigma^2) = (2\pi\sigma^2)^{-n/2} \exp\left\{ -\frac{1}{2\sigma^2} (\mathbf{Y}-X\boldsymbol{\beta})^T (\mathbf{Y}-X\boldsymbol{\beta}) \right\}. \]

Ignorando constantes:
\[ L(\boldsymbol{\beta},\sigma^2) \propto (\sigma^2)^{-n/2} \exp \left( -\frac{(\mathbf{Y}-X\boldsymbol{\beta})^T (\mathbf{Y}-X\boldsymbol{\beta})}{2\sigma^2} \right). \]

\vspace{0.4cm}

\subsection*{2. Especificação Bayesiana}
\addcontentsline{toc}{subsection}{2. Especificação Bayesiana}

Escolhemos o prior conjugado:
\[ \sigma^2 \sim IG(a_0,b_0), \] $a_0>0$ e $b_0>0$;

com densidade:
\[ p(\sigma^2) = \frac{b_0^{a_0}}{\Gamma(a_0)} (\sigma^2)^{-a_0-1} \exp\left( -\frac{b_0}{\sigma^2} \right). \]

Condicionalmente a $\sigma^2$:

\[
\boldsymbol{\beta}\mid\sigma^2
\sim
\mathcal{N}_p
\left(
\boldsymbol{\mu}_0,
\sigma^2\Lambda_0^{-1}
\right),
\]

onde $\Lambda_0\in\mathbb{R}^{p\times p}$ é uma matriz simétrica
\[
\Lambda_0=\Lambda_0^{\top}
\]
e definida positiva, isto é,

\[
\mathbf{x}^{\top}\Lambda_0\mathbf{x}>0,
\qquad
\forall\,\mathbf{x}\neq\mathbf{0}.
\]

Logo:
\[ p(\boldsymbol{\beta}\mid\sigma^2) = \frac{|\Lambda_0|^{1/2}}{(2\pi\sigma^2)^{p/2}} \exp \left[ -\frac{(\boldsymbol{\beta}-\boldsymbol{\mu}_0)^T \Lambda_0 (\boldsymbol{\beta}-\boldsymbol{\mu}_0)}{2\sigma^2} \right]. \]

Assim,
\[ (\boldsymbol{\beta},\sigma^2) \sim NIG(\boldsymbol{\mu}_0,\Lambda_0,a_0,b_0). \]

\vspace{0.4cm}

\subsection*{3. Aplicação do Teorema de Bayes}
\addcontentsline{toc}{subsection}{3. Aplicação do Teorema de Bayes}

Pelo Teorema de Bayes:

\[
p(\boldsymbol{\beta},\sigma^2\mid\mathbf Y)
=
\frac{
p(\mathbf Y\mid\boldsymbol{\beta},\sigma^2)
p(\boldsymbol{\beta},\sigma^2)
}{
p(\mathbf Y)
}
\]

Como \(p(\mathbf Y)\) não depende de
\(\boldsymbol{\beta}\) nem de \(\sigma^2\), usamos proporcionalidade:

\[
p(\boldsymbol{\beta},\sigma^2\mid\mathbf Y)
\propto
p(\mathbf Y\mid\boldsymbol{\beta},\sigma^2)
p(\boldsymbol{\beta},\sigma^2)
\]

Agora expandimos cada termo.

\vspace{0.3cm}

\textbf{Passo 1 — Escrever a verossimilhança}

Do modelo:

\[
\mathbf Y\mid\boldsymbol\beta,\sigma^2
\sim
\mathcal N_n
(X\boldsymbol\beta,\sigma^2I)
\]

a densidade é:

\[
p(\mathbf Y\mid\boldsymbol\beta,\sigma^2)
=
\frac1{(2\pi\sigma^2)^{n/2}}
\exp
\left(
-
\frac{
(\mathbf Y-X\boldsymbol\beta)^T
(\mathbf Y-X\boldsymbol\beta)
}{
2\sigma^2
}
\right)
\]

Ignorando constantes:

\[
p(\mathbf Y\mid\boldsymbol\beta,\sigma^2)
\propto
(\sigma^2)^{-n/2}
\exp
\left(
-
\frac{
(\mathbf Y-X\boldsymbol\beta)^T
(\mathbf Y-X\boldsymbol\beta)
}{
2\sigma^2
}
\right)
\]

\vspace{0.3cm}

\textbf{Passo 2 — Escrever o prior conjunto}

Como:

\[
p(\boldsymbol\beta,\sigma^2)
=
p(\boldsymbol\beta\mid\sigma^2)
p(\sigma^2)
\]

escrevemos cada densidade.

Para a variância:

\[
\sigma^2\sim IG(a_0,b_0)
\]

temos:

\[
p(\sigma^2)
=
\frac{b_0^{a_0}}
{\Gamma(a_0)}
(\sigma^2)^{-a_0-1}
\exp
\left(
-\frac{b_0}{\sigma^2}
\right)
\]

Ignorando constantes:

\[
p(\sigma^2)
\propto
(\sigma^2)^{-a_0-1}
\exp
\left(
-\frac{b_0}{\sigma^2}
\right)
\]

Agora:

\[
\boldsymbol\beta\mid\sigma^2
\sim
\mathcal N_p
(
\mu_0,
\sigma^2\Lambda_0^{-1}
)
\]

logo:

\[
p(\boldsymbol\beta\mid\sigma^2)
=
\frac{
|\Lambda_0|^{1/2}
}{
(2\pi\sigma^2)^{p/2}
}
\exp
\left(
-
\frac{
(\boldsymbol\beta-\mu_0)^T
\Lambda_0
(\boldsymbol\beta-\mu_0)
}{
2\sigma^2
}
\right)
\]

Ignorando constantes:

\[
p(\boldsymbol\beta\mid\sigma^2)
\propto
(\sigma^2)^{-p/2}
\exp
\left(
-
\frac{
(\boldsymbol\beta-\mu_0)^T
\Lambda_0
(\boldsymbol\beta-\mu_0)
}{
2\sigma^2
}
\right)
\]

Portanto:

\[
p(\boldsymbol\beta,\sigma^2)
\propto
(\sigma^2)^{-a_0-1-p/2}
\exp
\left[
-
\frac1{\sigma^2}
\left(
b_0+
\frac12
(\boldsymbol\beta-\mu_0)^T
\Lambda_0
(\boldsymbol\beta-\mu_0)
\right)
\right]
\]

\vspace{0.3cm}

\textbf{Passo 3 — Multiplicar prior × verossimilhança}

Aplicando:

\[
p(\boldsymbol\beta,\sigma^2\mid Y)
\propto
p(Y\mid\boldsymbol\beta,\sigma^2)
p(\boldsymbol\beta,\sigma^2)
\]

substituímos:

\[
\propto
(\sigma^2)^{-n/2}
\exp
\left(
-
\frac{
(Y-X\beta)^T(Y-X\beta)
}{
2\sigma^2
}
\right)
\]

\[
\times
(\sigma^2)^{-a_0-1-p/2}
\exp
\left[
-
\frac1{\sigma^2}
\left(
b_0+
\frac12
(\beta-\mu_0)^T
\Lambda_0
(\beta-\mu_0)
\right)
\right]
\]

Agrupando potências:

\[
(\sigma^2)^{-n/2}
(\sigma^2)^{-a_0-1-p/2}
=
(\sigma^2)^{-a_0-1-n/2-p/2}
\]

obtemos:

\[
p(\beta,\sigma^2\mid Y)
\propto
(\sigma^2)^{-a_0-1-n/2-p/2}
\]

\[
\times
\exp
\left[
-
\frac1{\sigma^2}
\left(
b_0
+
\frac12
(Y-X\beta)^T(Y-X\beta)
+
\frac12
(\beta-\mu_0)^T
\Lambda_0
(\beta-\mu_0)
\right)
\right]
\]

que é exatamente a forma desejada.

\vspace{0.4cm}

\subsection*{4. Completando o Quadrado}
\addcontentsline{toc}{subsection}{4. Completando o Quadrado}

Partimos do termo que aparece dentro da exponencial da posterior:

\[
(\mathbf Y-X\boldsymbol\beta)^T
(\mathbf Y-X\boldsymbol\beta)
+
(\boldsymbol\beta-\boldsymbol\mu_0)^T
\Lambda_0
(\boldsymbol\beta-\boldsymbol\mu_0)
\]

Nosso objetivo é reescrever essa expressão em função de

\[
(\boldsymbol\beta-\boldsymbol\mu_n)^T
\Lambda_n
(\boldsymbol\beta-\boldsymbol\mu_n).
\]

\vspace{0.3cm}

\textbf{Passo 1 — Expandir o primeiro termo}

Expandimos:

\[
(\mathbf Y-X\boldsymbol\beta)^T
(\mathbf Y-X\boldsymbol\beta)
\]

Distribuindo:

\[
=
(\mathbf Y^T-\boldsymbol\beta^TX^T)
(\mathbf Y-X\boldsymbol\beta)
\]

Multiplicando termo a termo:

\[
=
\mathbf Y^T\mathbf Y
-
\mathbf Y^TX\boldsymbol\beta
-
\boldsymbol\beta^TX^T\mathbf Y
+
\boldsymbol\beta^TX^TX\boldsymbol\beta
\]

Como:

\[
\mathbf Y^TX\boldsymbol\beta
=
\boldsymbol\beta^TX^T\mathbf Y
\]

(observação: ambos são escalares),

obtemos:

\[
(\mathbf Y-X\boldsymbol\beta)^T
(\mathbf Y-X\boldsymbol\beta)
=
\mathbf Y^T\mathbf Y
-
2\boldsymbol\beta^TX^T\mathbf Y
+
\boldsymbol\beta^TX^TX\boldsymbol\beta
\]

\vspace{0.4cm}

\textbf{Passo 2 — Expandir o segundo termo}

Agora expandimos:

\[
(\boldsymbol\beta-\boldsymbol\mu_0)^T
\Lambda_0
(\boldsymbol\beta-\boldsymbol\mu_0)
\]

Distribuindo:

\[
=
(\boldsymbol\beta^T-\boldsymbol\mu_0^T)
\Lambda_0
(\boldsymbol\beta-\boldsymbol\mu_0)
\]

Multiplicando:

\[
=
\boldsymbol\beta^T\Lambda_0\boldsymbol\beta
-
\boldsymbol\beta^T\Lambda_0\boldsymbol\mu_0
-
\boldsymbol\mu_0^T\Lambda_0\boldsymbol\beta
+
\boldsymbol\mu_0^T\Lambda_0\boldsymbol\mu_0
\]

Como:

\[
\Lambda_0=\Lambda_0^T
\]

segue que:

\[
\boldsymbol\mu_0^T\Lambda_0\boldsymbol\beta
=
\boldsymbol\beta^T\Lambda_0\boldsymbol\mu_0
\]

Logo:

\[
(\boldsymbol\beta-\boldsymbol\mu_0)^T
\Lambda_0
(\boldsymbol\beta-\boldsymbol\mu_0)
=
\boldsymbol\beta^T\Lambda_0\boldsymbol\beta
-
2\boldsymbol\beta^T\Lambda_0\boldsymbol\mu_0
+
\boldsymbol\mu_0^T\Lambda_0\boldsymbol\mu_0
\]

\vspace{0.4cm}

\textbf{Passo 3 — Somar os dois desenvolvimentos}

Somando:

\[
\mathbf Y^T\mathbf Y
-
2\boldsymbol\beta^TX^T\mathbf Y
+
\boldsymbol\beta^TX^TX\boldsymbol\beta
\]

\[
+
\boldsymbol\beta^T\Lambda_0\boldsymbol\beta
-
2\boldsymbol\beta^T\Lambda_0\boldsymbol\mu_0
+
\boldsymbol\mu_0^T\Lambda_0\boldsymbol\mu_0
\]

Agrupando os termos quadráticos:

\[
=
\boldsymbol\beta^T
(X^TX+\Lambda_0)
\boldsymbol\beta
\]

Agrupando os termos lineares:

\[
-
2\boldsymbol\beta^T
(X^T\mathbf Y+\Lambda_0\mu_0)
\]

Mantendo constantes:

\[
+
\mathbf Y^T\mathbf Y
+
\mu_0^T\Lambda_0\mu_0
\]

Portanto:

\[
=
\boldsymbol\beta^T
(X^TX+\Lambda_0)
\boldsymbol\beta
-
2\boldsymbol\beta^T
(X^T\mathbf Y+\Lambda_0\mu_0)
+
\mathbf Y^T\mathbf Y
+
\mu_0^T\Lambda_0\mu_0
\]

\vspace{0.4cm}

\textbf{Passo 4 — Definir os parâmetros posteriores}

Definimos:

\[
\boxed{
\Lambda_n=X^TX+\Lambda_0
}
\]

e

\[
\boxed{
\Lambda_n\mu_n
=
X^T\mathbf Y+\Lambda_0\mu_0
}
\]

equivalentemente:

\[
\boxed{
\mu_n
=
\Lambda_n^{-1}
(X^T\mathbf Y+\Lambda_0\mu_0)
}
\]

Substituindo:

\[
=
\boldsymbol\beta^T
\Lambda_n
\boldsymbol\beta
-
2\boldsymbol\beta^T
\Lambda_n\mu_n
+
\mathbf Y^T\mathbf Y
+
\mu_0^T\Lambda_0\mu_0
\]

\vspace{0.4cm}

\textbf{Passo 5 — Completar o quadrado}

Usamos a identidade matricial:

\[
(\beta-\mu)^TA(\beta-\mu)
=
\beta^TA\beta
-
2\beta^TA\mu
+
\mu^TA\mu
\]

Logo:

\[
\beta^T\Lambda_n\beta
-
2\beta^T\Lambda_n\mu_n
\]

\[
=
(\beta-\mu_n)^T
\Lambda_n
(\beta-\mu_n)
-
\mu_n^T\Lambda_n\mu_n
\]

Substituindo:

\[
=
(\beta-\mu_n)^T
\Lambda_n
(\beta-\mu_n)
+
\mathbf Y^T\mathbf Y
+
\mu_0^T\Lambda_0\mu_0
-
\mu_n^T\Lambda_n\mu_n
\]

que é exatamente o resultado desejado.

\vspace{0.4cm}

\subsection*{5. Atualização dos Hiperparâmetros}
\addcontentsline{toc}{subsection}{5. Atualização dos Hiperparâmetros}

Os parâmetros posteriores tornam-se:
\[ \boxed{a_n = a_0+\frac{n}{2}} \]
\[ \boxed{b_n = b_0 + \frac{1}{2} \left( \mathbf{Y}^T\mathbf{Y} + \boldsymbol{\mu}_0^T\Lambda_0\boldsymbol{\mu}_0 - \boldsymbol{\mu}_n^T\Lambda_n\boldsymbol{\mu}_n \right)} \]

\vspace{0.4cm}

\subsection*{6. Distribuição Posterior Conjunta}
\addcontentsline{toc}{subsection}{6. Distribuição Posterior Conjunta}

Concluímos:
\[ (\boldsymbol{\beta},\sigma^2\mid \mathbf{Y}) \sim NIG(\boldsymbol{\mu}_n,\Lambda_n,a_n,b_n). \]

Equivalente à fatoração:
\[ \sigma^2\mid \mathbf{Y} \sim IG(a_n,b_n) \]
\[ \boldsymbol{\beta}\mid\sigma^2,\mathbf{Y} \sim \mathcal{N}_p ( \boldsymbol{\mu}_n, \sigma^2\Lambda_n^{-1} ). \]

\vspace{0.4cm}
% \subsection{7. Distribuição Preditiva Posterior}

%Para nova observação:

%[x_\star^T\beta+\varepsilon_\star.]

%A distribuição preditiva é:

%\intp(Y_\star\mid\beta,\sigma^2)p(\beta,\sigma^2\mid Y)d\beta d\sigma^2.]

%Resultado:

%[\boxed{Y_\star\mid Y\simt_{2a_n}\left(x_\star^T\mu_n,\frac{b_n}{a_n}\left(1+x_\star^T\Lambda_n^{-1}x_\star\right)\right)}\]

\subsection*{7. Interpretação Bayesiana}
\addcontentsline{toc}{subsection}{7. Interpretação Bayesiana}

Os hiperparâmetros possuem interpretação direta:
\[ \boldsymbol{\mu}_0 \implies \text{média a priori dos coeficientes} \]
\[ \Lambda_0 \implies \text{precisão a priori} \]
\[ a_0, b_0 \implies \text{informação prévia sobre variância}. \]

A atualização:
\[ \boldsymbol{\mu}_n = (X^TX+\Lambda_0)^{-1} (X^T\mathbf{Y}+\Lambda_0\boldsymbol{\mu}_0) \]
mostra explicitamente que a posterior é uma média ponderada entre os \text{dados} e o \text{conhecimento prévio}.

No limite em que $\Lambda_0 \to 0$ (prior não-informativo), recupera-se a inferência clássica de mínimos quadrados ordinários.

\section*{Dados}
\addcontentsline{toc}{section}{Dados}

O conjunto contém $n = 9\,568$ observações e $p = 5$ variáveis (4 covariáveis
e 1 resposta), todas contínuas. Nenhum observação apresenta dados faltantes. 
 
A Tabela~\ref{tab:amostra} exibe as 10 primeiras e as 10 últimas observações do
conjunto de dados, na ordem original disponibilizada pelo UCI.
 
\begin{table}[ht]
\centering
\caption{Primeiras e últimas 10 observações do conjunto CCPP.}
\label{tab:amostra}
\renewcommand{\arraystretch}{1.25}
\begin{tabular}{cS[table-format=2.2]S[table-format=2.2]
                  S[table-format=4.2]S[table-format=3.2]S[table-format=3.2]}
  \toprule
  \textbf{Obs.} & \textbf{AT (°C)} & \textbf{V (cm Hg)}
                & \textbf{AP (mbar)} & \textbf{RH (\%)} & \textbf{PE (MW)} \\
  \midrule
  % ── Primeiras 10 ──────────────────────────────────────────────────────────
  1    & 14.96 & 41.76 & 1024.07 & 73.17 & 463.26 \\
  2    & 25.18 & 62.96 & 1020.04 & 59.08 & 444.37 \\
  3    &  5.11 & 39.40 & 1012.16 & 92.14 & 488.56 \\
  4    & 20.86 & 57.32 & 1010.24 & 76.64 & 446.48 \\
  5    & 10.82 & 37.50 & 1009.23 & 96.62 & 473.90 \\
  6    & 26.27 & 59.44 & 1012.23 & 58.77 & 443.67 \\
  7    & 15.89 & 43.96 & 1014.02 & 75.24 & 467.35 \\
  8    &  9.48 & 44.71 & 1019.12 & 66.43 & 478.42 \\
  9    & 14.64 & 45.00 & 1021.78 & 41.25 & 475.98 \\
  10   & 11.74 & 43.56 & 1015.14 & 70.72 & 477.50 \\
  \midrule[\heavyrulewidth]
  \multicolumn{6}{c}{\textit{(...)}} \\
  \midrule[\heavyrulewidth]
  % ── Últimas 10 ────────────────────────────────────────────────────────────
  9559 & 28.37 & 69.70 & 1007.34 & 77.60 & 432.87 \\
  9560 & 21.57 & 58.39 & 1016.26 & 77.60 & 445.11 \\
  9561 & 12.07 & 45.86 & 1014.19 & 78.43 & 468.49 \\
  9562 & 22.31 & 60.83 & 1013.20 & 72.60 & 444.82 \\
  9563 & 23.16 & 65.50 & 1008.74 & 60.05 & 441.81 \\
  9564 & 18.67 & 52.88 & 1005.71 & 89.25 & 453.68 \\
  9565 & 24.68 & 67.63 & 1010.85 & 72.44 & 435.16 \\
  9566 & 23.93 & 63.65 & 1011.41 & 71.38 & 441.17 \\
  9567 & 21.01 & 56.63 & 1014.68 & 71.72 & 448.14 \\
  9568 & 22.39 & 68.05 & 1008.08 & 81.23 & 436.54 \\
  \bottomrule
\end{tabular}
\end{table}

\section*{Análise Descritiva}
\addcontentsline{toc}{section}{Análise Descritiva}

\section*{Modelo de Regressão Linear}
\addcontentsline{toc}{section}{Modelo de Regressão Linear}

\subsection*{Notação e Definições}
\addcontentsline{toc}{subsection}{Notação e Definições}
 
Seja $i = 1, 2, \ldots, n$ o índice da observação, com $n = 9\,568$.
Definimos as seguintes quantidades para a $i$-ésima observação:
 
\begin{align*}
  Y_i    &\;:\; \text{produção líquida de energia (MW)}, \\
  x_{i1} &\;:\; \text{temperatura ambiente } AT_i \text{ (°C)}, \\
  x_{i2} &\;:\; \text{vácuo de exaustão } V_i \text{ (cm Hg)}, \\
  x_{i3} &\;:\; \text{pressão ambiente } AP_i \text{ (mbar)}, \\
  x_{i4} &\;:\; \text{umidade relativa } RH_i \text{ (\%)}.
\end{align*}
 
\subsection*{Especificação Escalar}
\addcontentsline{toc}{subsection}{Especificação Escalar}
 
 
O modelo de regressão linear múltipla para a $i$-ésima observação é:
\begin{equation}
  Y_i \;=\; \beta_0 + \beta_1 x_{i1} + \beta_2 x_{i2}
           + \beta_3 x_{i3} + \beta_4 x_{i4} + \varepsilon_i,
  \qquad i = 1, \ldots, n,
  \label{eq:modelo_escalar}
\end{equation}
onde:
 
\begin{itemize}
  \item $\beta_0 \in \mathbb{R}$ é o \textbf{intercepto} --- produção esperada
        quando todas as covariáveis são zero;
  \item $\beta_1 \in \mathbb{R}$ é o coeficiente de AT --- variação esperada em
        PE para cada aumento de $1\,\text{°C}$ em temperatura;
  \item $\beta_2 \in \mathbb{R}$ é o coeficiente de V --- variação esperada em
        PE para cada aumento de $1\,\text{cm\,Hg}$ no vácuo;
  \item $\beta_3 \in \mathbb{R}$ é o coeficiente de AP --- variação esperada em
        PE para cada aumento de $1\,\text{mbar}$ na pressão;
  \item $\beta_4 \in \mathbb{R}$ é o coeficiente de RH --- variação esperada em
        PE para cada aumento de $1\,\text{pp}$ na umidade relativa;
  \item $\varepsilon_i$ é o \textbf{erro aleatório} da $i$-ésima observação.
\end{itemize}
 
\subsection*{Especificação Matricial}
\addcontentsline{toc}{subsection}{Especificação Matricial}
 
Definindo os vetores e matrizes:
\begin{align*}
  \mY_{n \times 1}   &= (Y_1, Y_2, \ldots, Y_n)\T, \\[4pt]
  X_{n \times p}     &= \begin{pmatrix}
                          1 & x_{11} & x_{12} & x_{13} & x_{14} \\
                          1 & x_{21} & x_{22} & x_{23} & x_{24} \\
                          \vdots & \vdots & \vdots & \vdots & \vdots \\
                          1 & x_{n1} & x_{n2} & x_{n3} & x_{n4}
                        \end{pmatrix}, \\[4pt]
  \mb_{p \times 1}   &= (\beta_0, \beta_1, \beta_2, \beta_3, \beta_4)\T, \\[4pt]
  \beps_{n \times 1} &= (\varepsilon_1, \varepsilon_2, \ldots, \varepsilon_n)\T,
\end{align*}
com $p = 5$ (intercepto + 4 covariáveis), o modelo~\eqref{eq:modelo_escalar}
escreve-se compactamente como:
\begin{equation}
  \mY \;=\; X\mb + \beps.
  \label{eq:modelo_matricial}
\end{equation}
 
\subsection*{Distribuição dos Erros e Verossimilhança}
\addcontentsline{toc}{subsection}{Distribuição dos Erros e Verossimilhança}
 
Assumindo erros independentes e normalmente distribuídos:
\[
  \beps \sim \Normal_{n}(\mathbf{0}_n,\, \sig I_n),
  \qquad \sig > 0,
\]
a distribuição condicional de $\mY$ é:
\begin{equation}
  \mY \mid \mb, \sig, X \;\sim\; \Normal_{n}(X\mb,\, \sig I_n),
  \label{eq:likelihood_dist}
\end{equation}
com verossimilhança:
\begin{equation}
  L(\mb, \sig) \;\propto\;
  (\sig)^{-n/2}
  \exp\!\left(
    -\frac{(\mY - X\mb)\T(\mY - X\mb)}{2\sig}
  \right).
  \label{eq:vero}
\end{equation}

\newpage
\subsection*{Interpretação dos Coeficientes}
\addcontentsline{toc}{subsection}{Interpretação dos Coeficientes}
 
Com base nas correlações observadas na literatura para este conjunto de dados,
espera-se, a posteriori, que:
 
 
\begin{modelbox}
\textbf{Modelo completo:}
\[
  PE_i \;=\; \beta_0 + \beta_1\,AT_i + \beta_2\,V_i
            + \beta_3\,AP_i + \beta_4\,RH_i + \varepsilon_i,
\]
\[
  \varepsilon_i \overset{\mathrm{iid}}{\sim} \Normal(0, \sig),
  \qquad i = 1, \ldots, 9568.
\]
\textbf{Parâmetros:}
$\boldsymbol{\theta} = (\beta_0, \beta_1, \beta_2, \beta_3, \beta_4, \sig) \in \mathbb{R}^4 \times \mathbb{R}_{>0}^2$.
 
\medskip
\textbf{Dimensões:} $n = 9568$, $p = 5$
(1 intercepto $+ 4$ covariáveis), posto$(X) = p$
(assumido; verificável pela ausência de multicolinearidade perfeita).
\end{modelbox}
 
% ══════════════════════════════════════════════════════════════════════════════
\section*{Especificação Bayesiana}
\addcontentsline{toc}{section}{Especificação Bayesiana}
 
Adotando priors conjugados independentes:
\begin{align}
  \sig   &\;\sim\; \GI(a_0,\, b_0),
  \label{eq:prior_sig} \\[4pt]
  \mb \mid \sig &\;\sim\; \Normal_p\!\left(
                    \boldsymbol{\mu}_0,\;
                    \sig\,\Lambda_0^{-1}
                  \right),
  \label{eq:prior_beta}
\end{align}
de forma que $(\mb, \sig) \sim \mathcal{NIG}(\boldsymbol{\mu}_0, \Lambda_0, a_0, b_0)$.
 
A posterior conjunta tem forma fechada:
\[
  (\mb, \sig \mid \mY, X)
  \;\sim\;
  \mathcal{NIG}(\boldsymbol{\mu}_n, \Lambda_n, a_n, b_n),
\]
com hiperparâmetros atualizados:
\begin{align}
  \Lambda_n      &= \Lambda_0 + X\T X, \\
  \boldsymbol{\mu}_n &= \Lambda_n^{-1}
                        (\Lambda_0\boldsymbol{\mu}_0 + X\T\mY), \\
  a_n            &= a_0 + \frac{n}{2}, \\
  b_n            &= b_0 + \frac{1}{2}
                    \!\left(
                      \mY\T\mY
                      + \boldsymbol{\mu}_0\T\Lambda_0\boldsymbol{\mu}_0
                      - \boldsymbol{\mu}_n\T\Lambda_n\boldsymbol{\mu}_n
                    \right).
\end{align}
 
 
\subsection*{Escolha dos Hiperparâmetros}
\addcontentsline{toc}{subsection}{Escolha dos Hiperparâmetros}
 
Para análise inicial com pouca informação prévia sobre os coeficientes, uma
escolha razoável é:
\[
  \boldsymbol{\mu}_0 = \mathbf{0}_5,
  \quad
  \Lambda_0 = ,
  \quad
  a_0 = b_0 = 0{,}01.
\]

 





\end{document}
